\documentclass[11pt,letterpaper]{article}
\usepackage[margin=1.02in,top=0.91in,bottom=0.93in]{geometry}
\usepackage{lmodern}\pdfmapfile{+lm.map}
\usepackage{amsmath,amssymb,amsthm,booktabs,array,microtype,xcolor}
\usepackage[colorlinks=true,linkcolor=teal,urlcolor=teal,citecolor=teal]{hyperref}
\usepackage{fancyhdr}
\definecolor{deepblue}{HTML}{123E53}
\pagestyle{fancy}\fancyhf{}\renewcommand{\headrulewidth}{0.3pt}
\lhead{\small\color{deepblue}Methods from a formal proof of Dirichlet GRH}
\rhead{\small Research note}\cfoot{\small\thepage}
\setlength{\headheight}{14pt}
\setlength{\emergencystretch}{1.1em}
\setlength{\parindent}{1.15em}
\setlength{\parskip}{2pt}
\setcounter{tocdepth}{1}
\newtheorem{theorem}{Theorem}[section]
\newtheorem{proposition}[theorem]{Proposition}
\newtheorem{corollary}[theorem]{Corollary}
\newcommand{\src}[2]{\href{https://github.com/lsernau/grh-lean/blob/72fead102f663d5c7a407964a13dc0cf9d977207/grh/src/#1}{#2}}
\newcommand{\repo}[2]{\href{https://github.com/lsernau/grh-lean/blob/72fead102f663d5c7a407964a13dc0cf9d977207/#1}{#2}}
\newcommand{\incoming}[2]{\href{https://github.com/VladimirReshetnikov/ProveIt/blob/5cc267af75a635199ac5d93a7d49cbce521daefc/docs/incoming/#1}{#2}}
\hypersetup{pdftitle={Beyond Dirichlet GRH: zero multiplicities, cyclotomic annihilators, and a prime-power race},pdfauthor={Research analysis prepared for Vladimir Reshetnikov}}
\begin{document}
\thispagestyle{plain}
\noindent{\small\sffamily\color{deepblue}RESEARCH NOTE\hfill 26 NOVEMBER 2026}\par
\vspace{0.5cm}
\noindent{\huge\bfseries Beyond Dirichlet GRH}\par\vspace{4pt}
\noindent{\Large Zero multiplicities, cyclotomic annihilators,\\[2pt]
and a prime-power race}\par\vspace{0.32cm}
\noindent Prepared for Vladimir Reshetnikov\par
\vspace{0.25cm}\hrule\vspace{0.25cm}
\noindent\textbf{Abstract.}
For Galois extensions with group $S_n$, $3\le n\le12$, explicit positive
character identities prove the holomorphy away from $1$ of
$\zeta_E^{\,n-1}\zeta_F/(\zeta_L\zeta_K)$, where $E$ is the fixed field
of a transposition and $F$ the quadratic subfield. This resolves the cases $3\le n\le12$
of the earlier ProveIt question on the optimal cost of detecting zeros.
Inside the upstream cyclotomic argument we determine the exact gcd of a
family of Stickelberger witnesses. At $(N,e)=(5,1)$ the least individual
witness is $8400$, whereas their gcd is $1200$. For Suzuki's unsmoothed
weighted prime sum we prove that its error is negative with lower logarithmic
density exceeding $99.992057\%$, using finite resonance checks on zeta zeros
and no global linear-independence hypothesis. A combinatorial appendix
verifies finite higher-rank cases of Kova\v cevi\'c's anticode conjecture
with all equality cases, and determines the minimum order of an abelian group admitting a five-element
$B_3$ set.
\tableofcontents
\clearpage
\section{Where the Lean proof enters}
\label{sec:source}
The readable edition's \src{Statements.lean}{statement file} imports Mathlib.
Its \src{Main.lean}{main file} exports five results: Mathlib's RH,
the equivalent explicit zeta-zero statement, nonvanishing to the right
of $1/2$ for Dirichlet characters away from the principal pole, the
open-strip zeros-on-the-line statement for all such characters, and its
primitive specialization. All positive conductors or moduli are allowed.
The \repo{grh/README.md}{edition README} reports the three axioms
\texttt{propext}, \texttt{Classical.choice}, and \texttt{Quot.sound}, and
reports independent proof-term checks. We inspected the source at
\texttt{72fead102f66}. The upstream build and independent kernel replays
are the maintainers' checks. The new arguments here are mathematical proofs
with accompanying computations, outside that Lean build.

\paragraph{From heights to zero multiplicities.}
The geometric part compares upper and lower bounds for the covolume of
one lattice of automorphic forms. Chains of theta lifts, local
divisibility of derivatives and a second-moment spanning argument supply
the short forms; the determinant estimate for integral forms supplies
the opposing bound. Canceling a common main term leaves a one-sided
estimate for a smoothed von Mangoldt sum. A chain of length $m$ with
$\sqrt x\le m\le2\sqrt x$ gives the decisive scale. In the analytic
step a zero becomes a pole of a transform with residue proportional to
its multiplicity. A one-sided bound and positivity at real poles exclude
zeros in the desired half-plane. See the
\repo{grh/GUIDE.md}{reading guide},
\src{New/Reduction/MainInequality.lean}{main inequality},
\src{New/Reduction/OneSidedBound.lean}{prime bound},
and \src{New/Reduction/Zeros/Criterion.lean}{zero criterion}.

\paragraph{What transfers in each argument.}
The reflection result uses zero multiplicities as local data, then finds
positive combinations of one-dimensional subgroup characters that retain
them. The cyclic character of the quadratic subfield is exactly where
Dirichlet GRH enters its off-line corollary. The cyclotomic calculation
uses a different part of the formalization, isolating the two quantifiers
in the finite-level integer used in the Kolyvagin argument. The prime-race
estimate uses the zeta instance of the final theorem, Suzuki's explicit
formula, and verified moments. The appendix takes inspiration from the
bounded-coefficient modular encodings in the proof's Behrend construction;
its upper bounds come from an independent lattice argument.

\paragraph{Comparison set.}
We inspected the 46 ZIP entries in the
\href{https://github.com/VladimirReshetnikov/ProveIt/tree/5cc267af75a635199ac5d93a7d49cbce521daefc/docs/incoming}{incoming snapshot at \texttt{5cc267af75a6}}.
Two specific baselines matter. The
\incoming{grh_reflection_transfer_2026-11-25.zip}{reflection report}
asks for the optimal factor $n-1$. The
\incoming{grh_cm_four_terms_and_prime_bias.zip}{prime-bias report}
already proves that both signs have positive lower logarithmic frequency
and bounds the rare sign by $0.00323$. Earlier discussions of Vandiver
identify a parity and comparison problem. Sharpening an annihilator of the
specified odd cocycles is meaningful independently of that problem.
We use published antecedents and this pinned comparison set in stating
the advances; global priority has not been established.
\clearpage
\section{Sharp zero multiplicities for symmetric Galois groups}
\label{sec:reflection}

An incoming report considers the splitting field of a polynomial with
symmetric Galois group. It asks whether the order of an off-line zero
can be recovered with the optimal loss from one much smaller set of
Frobenius classes. Its symmetric-group estimate is
$2(n-H_n)$, where $H_n=\sum_{j=1}^n1/j$.
It gives $n-1$ as the best possible lower bound for a proof from local
character positivity. The following result attains that bound for
every $3\le n\le 12$ and also controls zeros on the critical line. The precise
antecedent is \emph{Reflection detectors for nonabelian Riemann
hypotheses}, the symmetric-group corollary and the question in its final section,
in the \href{https://github.com/VladimirReshetnikov/ProveIt/blob/5cc267af75a635199ac5d93a7d49cbce521daefc/docs/incoming/grh_reflection_transfer_2026-11-25.zip}{pinned reflection report}.

\paragraph{Theorem.}
Let $L/K$ be a finite Galois extension of number fields with group
$G=S_n$, $3\le n\le 12$. Fix a transposition $t$, set
$E=L^{\langle t\rangle}$ and $F=L^{A_n}$, and let $\varepsilon$
be the sign character. Then the meromorphic function
\[
 \mathcal Q(s)=\frac{\zeta_E(s)^{\,n-1}L(s,\varepsilon,L/K)}
                         {\zeta_L(s)}
 =\frac{\zeta_E(s)^{\,n-1}\zeta_F(s)}{\zeta_L(s)\zeta_K(s)}
\]
is holomorphic on $\mathbb C\setminus\{1\}$, with a pole of exact order
$n-2$ at~1. In particular, for $K=\mathbb Q$ and $1/2<\Re\rho<1$,
Dirichlet GRH gives
\[
 \operatorname{ord}_\rho\zeta_L\le (n-1)\operatorname{ord}_\rho\zeta_E .
\]
For example, the constants are $4$ and $5$ for $S_5$ and $S_6$. Each constant $n-1$ is optimal among deductions
using positivity against all characters induced from one-dimensional
subgroup characters and vanishing of the two global linear coefficients.

\paragraph{Arithmetic input.}
For $s\ne1$ let $\Theta=\sum_\chi (\operatorname{ord}_{s}L(\cdot,\chi,L/K))\chi$,
the sum over irreducible characters of $G$.
Brauer induction gives meromorphy, so this virtual character exists.
For every $H\le G$ and every linear character $\eta$ of $H$, Artin
induction formalism and Hecke holomorphy away from 1 give
$\langle\Theta|_H,\eta\rangle\ge0$. The same holds for
$\Psi=\varepsilon\Theta$ since $\varepsilon|_H\eta$ is linear.
The regular and permutation characters give
\[
 m:=\Psi(1)=\operatorname{ord}_s\zeta_L,\qquad
 a:=\tfrac12(m-\Psi(t))=\operatorname{ord}_s\zeta_E,\qquad
 b:=\langle\Psi,1\rangle=\operatorname{ord}_s L(\cdot,\varepsilon,L/K).
\]
These are the classical Heilbronn-character identities.
See H.~Heilbronn, \href{https://doi.org/10.4153/CJM-1973-090-3}{\emph{On real
zeros of Dedekind $\zeta$-functions}}, Canad. J. Math. 25 (1973), 870--873,
and R.~Foote, H.~Ginsberg and V.~K.~Murty,
\href{https://doi.org/10.1090/BULL/1492}{\emph{Heilbronn characters}}.

\paragraph{Finite proof.}
For a conjugacy type $T$, let $V_T$ be the value of $\Psi$ on $T$
and $c_T$ the size of the class. Write $\delta_T$ for its coordinate
unit vector. Given a subgroup $H$ and a surjection $e:H\to\mathbb Z/d$,
put
\[
 B_{H,e}(T)=\sum_{h\in H\cap T}\frac1{\varphi(d)}
       \sum_{u\in(\mathbb Z/d)^\times}
                   \exp(2\pi i u e(h)/d).
\]
The row $B_{H,e}$ is real and $B_{H,e}\cdot V\ge0$ by the preceding
local positivity. The sum over $u$ includes the complex conjugate of
each character. We claim the following exact vector identity in each certified degree. The table gives the positive weights for the two smallest nonsolvable groups.
\[
 \frac{n-3}{n-1}\delta_{1}-\delta_t+
             \frac{2}{(n-1)n!}(c_T)_T
       =\sum_{(H,e)} w_{H,e} B_{H,e}.                         \tag{R}
\]
Pairing with $V$ gives $2(a+(b-m)/(n-1))\ge0$, i.e.
$m\le(n-1)a+b$ at every $s\ne1$. This proves the holomorphy.
At $s=1$, each Dedekind zeta has a simple pole with nonzero residue,
which gives the stated exact order.

The notation $g_j$ in the following table means that the displayed
permutation maps to $j\bmod d$. Each row generates its indicated
subgroup and character. For example, the first row contains the identity,
two transpositions with character value $-1$, and one double
transposition with value $+1$.

\begin{center}\small
\begin{tabular}{r r r l}
\hline
$n$ & $w_{H,e}$ & $d$ & Generators with their images\\ \hline
5 & $3/8$&2 & $(45)_1,\ (23)_1$\\
5 & $1/16$&2 & $(45)_1,\ ((24)(35))_1$\\
5 & $1/24$&2 & $(45)_1,\ ((12)(34))_1$\\
5 & $1/40$&2 & $((23)(45))_1,\ ((12)(34))_1$\\ \hline
6 & $7/30$&2 & $(56)_1,\ ((23)(45))_1$\\
6 & $1/10$&3 & $(456)_2,\ ((12)(34)(56))_0$\\
6 & $1/45$&6 & $(456)_4,\ ((14)(25)(36))_3$\\
6 & $1/50$&2 & $((34)(56))_1,\ ((23)(45))_1$\\
6 & $1/5$&2 & $(56)_1,\ (34)_1,\ (12)_1$\\
6 & $1/40$&2 & $(56)_0,\ (34)_1,\ ((13)(24))_1$\\ \hline
\end{tabular}
\end{center}

For $S_5$ one can check (R) directly. In the type order
$1,2,3,22,4,32,5$ (fixed points suppressed), the four rows are
\[
\begin{pmatrix}
1&-2&0&1&0&0&0\\
1&-2&0&-1&2&0&0\\
1&-2&2&-3&0&2&0\\
1&0&0&-5&0&0&4
\end{pmatrix}.
\]
The class sizes are $(1,10,20,15,30,20,24)$.
To reproduce the checks for every degree, start with the identity
permutation at image zero and multiply by each generator, adding its
image modulo $d$, until the list stops growing. Each occurrence of an
already listed permutation has the same image. Sum the real cyclotomic
traces in each conjugacy class and check (R) entry by entry. The
accompanying \texttt{check\_reflection.py} performs these operations in
exact rational arithmetic, using the completely specified inputs in
\texttt{certificates.json}. The optimality witness is the usual standard
character $\Psi(g)=\#\operatorname{Fix}(g)-1$. For every $n\ge3$ it is
orthogonal to both linear characters, it satisfies every subgroup
positivity inequality, and it has $(m,a,b)=(n-1,1,0)$. This witnesses
sharpness of the group-theoretic certificate. Existence of an actual
off-line zero realizing equality is a separate question.

For $n=3$ the positive row is the sign character itself, weighted
by $1/6$ and aggregated by type. The accompanying data for each
remaining degree contain a list of generators and a cyclic character
on each subgroup in (R), as in the table. The number of columns in
(R) is the number of partitions of $n$, and its class sizes follow
from $c_T=n!/\prod_{j} j^{m_j}m_j!$ for a type with $m_j$ cycles
of length $j$. The checker enumerates the generated subgroup, with two class-count
shortcuts for large rows. A full $S_k$ is recognized by a $k$-cycle
and an adjacent transposition; its row uses the class-size formula and
the indicated sign power. An additional generator exchanging two
disjoint equal blocks gives $(S_k\times S_k)\rtimes C_2$. In the
base subgroup the cycle types concatenate. In the exchange coset the
cycles have twice the lengths of the cycles of the product of the
two component permutations, and there are $k!$ pairs with each
given product. The linear character is the product of its values on
the two components, times its value on the exchange in that coset.
The checker matches the supplied generators and their character values
before applying either identity. No numerical optimization enters its
acceptance test.

\begin{center}\small
\begin{tabular}{c|rrrrrrrrrr}
$n$&3&4&5&6&7&8&9&10&11&12\\
conjugacy classes&3&5&7&11&15&22&30&42&56&77\\
positive subgroup rows&1&3&4&6&10&18&26&38&48&71
\end{tabular}
\end{center}

\paragraph{The character behind the calculation.}
Write $\sigma$ for the nontrivial character of $\langle t\rangle$.
The character whose pairing with $\Psi$ is
$(n-1)a-m+b$ equals
\[
 U_n=(n-1)\operatorname{Ind}_{\langle t\rangle}^{S_n}\sigma
                   -\operatorname{Reg}_{S_n}+1.
\]
It is a genuine character for every $n$. Indeed, take the Cayley graph
of $S_n$ for the $n-1$ adjacent transpositions, with left multiplication
on vertices. Its vertex module is regular and its oriented edge module
has one summand $\operatorname{Ind}\sigma$ for each generator, because
the edge stabilizer reverses the edge. The graph is connected, so the
cellular chain sequence gives $U_n=H_1$ at the level of characters.
A proof that $U_n$ belongs to the positive rational cone of characters
induced from linear subgroup characters for all $n$ would give the
multiplicity theorem in all degrees. Genuine character status alone
gives no such decomposition. Formula (R) supplies it in the stated
degrees by Frobenius reciprocity.

There is an obstruction to deducing the required decomposition from graph
homology in general. Let $A_5$ fix two vertices and permute five parallel
edges in its natural action. Its first homology is the irreducible
four-dimensional augmentation representation $U$ (the action is doubly
transitive). If $U$ were a positive rational sum of monomial characters,
each summand would be a multiple $kU$. Frobenius reciprocity gives
$1\le k\le4$, so the inducing subgroup would have index $4k$.
Indices $8$ and $16$ do not divide $60$, and an index-four subgroup would
give a nontrivial map of the simple group $A_5$ into $S_4$. In the
remaining case the subgroup has order five and $k=3$. The restriction of
$U$ to this subgroup is its regular representation minus its trivial
character, so every linear character occurs at most once. This rules
out $k=3$. Hence the general-degree problem needs more about the
adjacent-transposition graph than connectedness and a group action.

\paragraph{What the improvement detects.}
For the $S_5$ splitting field of $X^5-X-1$ in the prior report, an
off-line zero of multiplicity $m$ would have multiplicity at least
$\lceil m/4\rceil$ in the degree-60 transposition fixed field.
The former constant was $163/30$. For $S_6$ the corresponding degree
is $360$ and the former constant was $71/10$.
Artin--Hecke formalism gives the new meromorphic quotient over
every base field, independently of GRH. Over $\mathbb Q$, the Lean
Dirichlet result removes its only possible linear-character correction
to the right of the critical line. In the previous report's proposed
one-sided prime estimate, this identifies exactly how much multiplicity
the selected Frobenius classes must retain. An estimate for their prime
sum is still required to deduce RH for the full splitting field.

\clearpage
\section{An exact refinement of the cyclotomic annihilator}
\label{sec:cyc}
An intermediate step of the readable edition constructs an integer $X$ from a
Bernoulli matrix and proves that $4X^2$ kills a family of Galois cocycles.
Here we solve an optimization problem inside this argument. The value of a
single integer meeting the Lean hypotheses and the effective integer obtained
by combining several such proofs can differ. At modulus five and weight one
we determine both exactly. The statements in this section use the cyclotomic
lemmas imported by the proof. They are independent of zero location in the
critical strip.

\subsection{The upstream interface}
Fix integers $N,e\ge1$. Put $U=(\mathbb Z/N\mathbb Z)^\times$ and use the
representatives $0\le t<N$. Write
\[
 b(t)=(e+2)!N^{e+1} B_{e+1}(t/N),\quad K=(e+2)!(e+1)N,\qquad
 F(c)=(c^e-1)(c^{e+1}-1).
\]
Here $B_j(x)$ is the Bernoulli polynomial with $B_1(x)=x-1/2$.
The values $b(t)$ are integers. Call a character relevant if
$\chi(-1)=(-1)^{e+1}$, and put
$B_\chi=\sum_{t\in U} b(t)\chi(t)$. These sums are nonzero for all relevant
characters by the functional equation and nonvanishing at $1+e$.
The principal character for $N=1$ uses $t=0$. Characters at other moduli
extend by zero on nonunits.

The source's predicate \texttt{IsStickelbergerBound} asks a positive
candidate $X$ to satisfy two requirements.
First, for every rational prime $p$, some $c\equiv1\pmod N$ coprime to $p$
has $F(c)\mid X$. Such a $c$ must be at least two for positive $X$.
Second, for each $c\equiv1\pmod N$ and each relevant character
in a characteristic-zero integral domain $O$, there must be $w\in O$
independent of the modulus $m$ such that
\begin{equation}
 (c^{e+1}-1)X\equiv 2w S_{m,c,\chi}\pmod{mO},\quad
 S_{m,c,\chi}=\sum_{0\le a<Nm}\left\lfloor\frac{ca}{Nm}\right\rfloor
 a^e\chi(a),\quad (m,c)=1.                         \label{eq:stick}
\end{equation}
The source proves
$Kc^eS_{m,c,\chi}\equiv(c^{e+1}-1)B_\chi\pmod{mO}$.
It chooses the matrix
\[
 T_{u,v}=2b(t)+(-1)^e\mathbf1_{t=1}+\mathbf1_{t=-1},\qquad
 t=vu^{-1}\pmod N ,
\]
and supplies the witness
$X_0=F(1+N)F(1+2N)|\det T|$. Character vectors diagonalize $T$.
Relevant characters have eigenvalue $2B_\chi$, and the others have
eigenvalue $2(-1)^e$. The adjugate clears every relevant eigenvalue at once.
See \src{Lib/Cyclotomic/Voronoi.lean}{Voronoi, lines 134--151, 245--256 and 384--459}.

\subsection{The common denominator of the character values}
For each complex relevant $\chi$, let $d_\chi$ be the order of its image
and let $D_\chi$ generate the positive ideal
$(2B_\chi)\mathbb Z[\zeta_{d_\chi}]\cap\mathbb Z$.
Set $D=\mathop{\rm lcm}_\chi D_\chi$, with empty lcm one.
Equivalently $D_\chi$ is the least common denominator of the rational
coefficients of $(2B_\chi)^{-1}$ in the cyclotomic power basis. This gives a
finite exact calculation for every $N,e$.

\begin{proposition}
$D$ is the least positive integer divisible by every relevant $2B_\chi$ in
every characteristic-zero character value domain. For any coprime $a,b\ge2$
congruent to $1\pmod N$,
\[
             X(a,b)=\mathop{\rm lcm}(D,F(a),F(b))
\]
satisfies both upstream requirements.
\end{proposition}
\begin{proof}
The finite image of any such character generates a cyclotomic subring of
$O$. Indeed, in its fraction field the image is cyclic, and a primitive
$d$th root in characteristic zero has minimal polynomial $\Phi_d$ over
$\mathbb Q$. Thus each complex-character divisibility transfers to $O$.
Taking $O=\mathbb Z[\zeta_{d_\chi}]$ proves the converse. An algebraic
integer in $\mathbb Q(\zeta_d)$ belongs to $\mathbb Z[\zeta_d]$, so one
may equally compute the denominators in a single cyclotomic field containing
all character values.
For $X=X(a,b)$ the quotient $u=X/(2B_\chi)$ lies in $O$.
Multiply the upstream congruence by $2u$ to obtain \eqref{eq:stick} with
$w=uKc^e$. Each rational prime is coprime to at least one of $a,b$,
giving the first requirement. The adjugate argument also shows $D\mid|\det T|$.
\end{proof}

\subsection{An exact local calculation}
Define $M=\prod_p p^{m_p}$ using the following finite prescription.
Write $v_p$ for the ordinary integer valuation. For odd $p$,
\begin{equation}
m_p=\begin{cases}
2v_p(N)+v_p(e(e+1)) & p\mid N,\\
\displaystyle\sum_{j\in\{e,e+1\},\ p-1\mid j}(1+v_p(j)) & p\nmid N.
\end{cases}                                                    \label{eq:mp}
\end{equation}
For $p=2$ put $m_2=v_2(e(e+1))+2v_2(N)$ if $4\mid N$, and
$m_2=v_2(e(e+1))+3$ otherwise. Set $H=\mathop{\rm lcm}(D,M)$.

\begin{theorem}\label{th:gcd}
The gcd of the integers $X(a,b)$ over all admissible coprime pairs is exactly
$H$. Every positive integer satisfying the upstream first requirement is a
multiple of $M$.
\end{theorem}
\begin{proof}
For each $p$ minimize $v_p(F(c))$ over $c\ge2$ with $c\equiv1\pmod N$
and $p\nmid c$. For odd $p\mid N$, the lifting-the-exponent formula gives
$2v_p(c-1)+v_p(e(e+1))$ and one can take $v_p(c-1)=v_p(N)$.
For odd $p\nmid N$, an exponent divisible by $p-1$ forces at least
$1+v_p(j)$ in $c^j-1$ by Fermat's theorem and lifting the exponent.
A primitive root modulo $p^2$, combined with the progression modulo $N$,
attains equality for all $j$ at once and contributes zero for the other
exponents. At two, let $a'=v_2(c-1), b'=v_2(c+1)$. The two exponents have
opposite parity, giving total valuation
$2a'+b'-1+v_2(e(e+1))$. If $4\mid N$ we can take
$a'=v_2(N), b'=1$. Otherwise $c\equiv3\pmod8$ is possible and minimizes
$2a'+b'$ at four. This proves the formulas and $M\mid X$ for every
positive $X$ passing the first requirement.

Fix $p$ and choose one minimizing $a$. Choose $b>1$ congruent to $a$ modulo
$p^{m_p+1}$ and to $1$ modulo $N$ and every prime divisor of $a$. These
congruences are compatible, including when $p\mid N$, and give
$(a,b)=1$ and $v_p(F(b))=m_p$. Thus
$v_p(X(a,b))=\max(v_p(D),m_p)$ is attained for each $p$.
Taking the gcd proves the claim. A finite set of witnesses always suffices:
start with any one witness and minimize at each of its prime divisors.
\end{proof}

\begin{corollary}\label{cor:kill}
In the setting of the repository's odd component with coefficients
$O/\ell^\nu$, the integer $4H^2$ annihilates all its cocycles. For
$\ell>N+e+3$, at most $v_\ell(H)+1$ Galois elements suffice to detect an
element by evaluation. For a CM subfield $E\subseteq\mathbb Q(\mu_N)$,
under the source's coefficient hypotheses, $16\varphi(N)H^2$ annihilates
the module \texttt{shaOdd} at weight $e$ and every level.
\end{corollary}
\begin{proof}
For each genuine witness $X$, the source proves the assertions with $X$ in
place of $H$. Take the gcd of the integer annihilators using B\'ezout's
identity. For detection choose a witness attaining the minimum valuation at
the specified $\ell$. These two operations are different. In particular
the proof never requires $H$ itself to pass the first requirement.
See \src{Lib/Cyclotomic/Kolyvagin/ReciprocityComponentBounds.lean}{component
annihilation and detection, lines 676 and 1104}, and
\src{Lib/Cyclotomic/OddSha.lean}{CM descent, line 535}.
\end{proof}
Here $O$ is the integer ring of a finite extension of $\mathbb Q_\ell$ as in
the source. The component consists of locally constant cocycles on
$G_\mathbb Q$ for $(\mathrm{cyc}^e\chi)^{-1}$, of odd parity, subject to its
specified vanishing at conjugation, inertia and the local set at $\ell$.
CM descent also asks for enough roots of unity in $O$. The finite-level
CM object uses cocycles with its particular local conditions.
The identification of this object with a conventional cohomology group,
and control of even class-group eigenspaces in Vandiver's conjecture,
would require further arguments. The upstream
\src{Lib/Cyclotomic/OddSha.lean}{module overview} states the comparison
issue explicitly. Classical higher Stickelberger theory already gives arithmetic annihilation results (see \href{https://archive.numdam.org/articles/10.1016/j.crma.2003.09.019/}{Banaszak, \emph{Higher analogues of Stickelberger's theorem}}). Our sharpness statement concerns the
displayed construction.

\subsection{The smallest witness at modulus five}
At $N=5,e=1$, the four $b(t)$, in increasing order of $t$, are
$1,-11,-11,1$. The trivial and quadratic characters give relevant
eigenvalues $-40$ and $48$. The two odd characters give $-2$ each.
Thus $D=240$, $|\det T|=7680$, $F(6)=175$, $F(11)=1200$ and
$X_0=1\,612\,800\,000$.

\begin{theorem}
At $(N,e)=(5,1)$ the least positive integer satisfying the full upstream
predicate is $8400$. The gcd of \emph{all} its positive witnesses is $1200$.
\end{theorem}
\begin{proof}
The pair $(6,11)$ gives $8400$ and $(11,16)$ gives $61200$ since
$F(16)=3825$. Their gcd is $1200=M$, a divisor of every witness by
Theorem~\ref{th:gcd}. For $0<X\le8400$, the only possible auxiliary
integers are $6,11,16$: $F$ is increasing and $F(21)=8800$.
The eligible integers with $F(c)\mid X$ must have gcd one. The only way
within the budget is to include $6$ and $11$: $6$ and $16$ share a
factor, and $\mathop{\rm lcm}(F(11),F(16))=61200$.
Thus $8400\mid X$.
\end{proof}
The example proves more than a reduction in the written size of a bound.
The original witness permits the exceptional prime seven. Since
$7\nmid4\cdot1200^2$, the corresponding odd components vanish at seven at
\emph{every} finite level. More generally at modulus five and weight one
only $2,3,5$ can support these components. This is an exclusion of primes;
it does not assert a nonzero component at any of those three.

\begin{center}
\begin{tabular}{rr r r r r}
\toprule
$N$ & $e$ & $D$ & $M$ & $X(1+N,1+2N)$ & $H$\\
\midrule
5&1&240&1200&8400&1200\\
5&2&1440&1200&41\,176\,800&7200\\
7&2&16\,128&2352&1\,986\,179\,328&112\,896\\
8&1&192&384&23\,040&384\\
12&2&17\,280&864&2\,973\,559\,680&17\,280\\
15&1&2880&3600&489\,600&14\,400\\
\bottomrule
\end{tabular}
\end{center}
The checker in \path{checks/cyclotomic.py} uses polynomial inversion modulo
cyclotomic polynomials over $\mathbb Q$. It exhibits finite lists of
coprime-pair witnesses giving each claimed gcd. In the classical scalar
case $(N,e)=(1,11)$ it gives
$H=2^{10}3^4\cdot5\cdot7\cdot11\cdot13\cdot691$.
Consequently at primes greater than $15$ only $691$ survives this test.
It comes from the numerator of $B_{12}=-691/2730$.
This recovers the familiar Bernoulli obstruction in this normalization.

\clearpage
\section{A fortyfold improvement for Suzuki's unsmoothed prime-power race}
\label{sec:suzukimoments}
Write $\Lambda$ for the von Mangoldt function and set
\[
 E(x)=\sum_{n\le x}'\frac{\Lambda(n)}{\sqrt n}-2\sqrt x .
\]
The prime gives half weight to a final prime power. Suzuki asks how often
$E(x)\le0$ in logarithmic measure (\href{https://arxiv.org/html/2411.07436v3}{\S1.1, after (6)}).
The incoming report \texttt{grh\_cm\_four\_terms\_and\_prime\_bias.zip}
proves an upper frequency of $0.00323$ for the opposite sign and a
strictly positive lower frequency. We give an explicit lower bound and reduce the upper bound by using
finite information about additive relations between zero ordinates.

\medskip\noindent\textbf{Theorem.}
Let $p(a,L)=L^{-1}|\{t\in[a,a+L]:E(e^t)>0\}|$ for $a\ge1$. Then
\[
 10^{-937}<\liminf_{L\to\infty}\ \inf_{a\ge1}p(a,L)
 \le\limsup_{L\to\infty}\ \sup_{a\ge1}p(a,L)<0.00007943 .
\]
In particular, the lower logarithmic density of the set $E(x)<0$ exceeds
$99.992057\%$. Its upper logarithmic density is at most $1-10^{-937}$. The existence and exact value of the logarithmic density
at zero remain open in this analysis. The only global zero-location input
is the zeta instance of the Lean theorem. All independence information
used below is finite and verified.

\paragraph{The mean and its Fourier series.}
Write $\gamma$ for a distinct positive ordinate and $m_\gamma$ for its
multiplicity. Suzuki's explicit formula (36) and central value (47) give
\[
 E(e^t)=-c+F(t)+o(1),\qquad
 c=\tfrac12\bigl(\pi/2+\log(8\pi)+\gamma_{\mathrm E}\bigr)>2.686,
 \qquad F(t)=-2\sum_{\gamma>0}\frac{m_\gamma}{\gamma}\sin(\gamma t).
\]
Here $\gamma_{\mathrm E}$ is Euler's constant. This uses the same passage
from zero multiplicities to prime-sum residues that occurs in
\texttt{New/Reduction/Zeros/PrimeError.lean}, with the weight from Suzuki's
question. For clarity, the frequency-series convergence needed below can
be checked directly. For a finite sum $P=\sum A_\lambda e^{i\lambda t}$
put $B_j=\sum_{j\le\lambda<j+1}|A_\lambda|$. Integrating $|P|^2$ against
translates of $(\sin(t/2)/(t/2))^2$, whose Fourier transform is supported
in $[-1,1]$, gives
\[
 \sup_u\int_u^{u+1}|P(t)|^2dt\ll\sum_j B_j^2 .
\]
The kernel is bounded below on $[-1/2,1/2]$ and only equal or neighboring
frequency blocks contribute. The zero count in a unit block is
$O(\log(2+|j|))$, with multiplicities. Hence the symmetric partial sums
converge in this uniform local mean-square norm, with their limit agreeing
almost everywhere with the pointwise explicit formula. Finite-polynomial
means on long intervals converge uniformly in the starting point, by
integrating each exponential. Equivalently, these means are Haar
integrals on the closure of the phase flow
$t\mapsto(e^{i\gamma t})_{\gamma>0}$ in the product of circles. Uniform
mean-square approximation extends this interpretation to the positive
part of any sine-series tail, since $x\mapsto x_+$ is $1$-Lipschitz.
Write $\|\cdot\|_p$ for the Haar norm. Inversion in the phase group
sends every sine sum $Y$ to $-Y$, giving
$\|(Y)_+\|_p=2^{-1/p}\|Y\|_p$ whenever the moment exists.

\paragraph{Moment estimates from finite certificates.}
Number the zeros by increasing positive ordinate. They are simple over
the entire finite range used below. Decompose
$F=P+M_1+M_2+M_3+R$. The indices for $P$ are 1 through 500. Each $M_j$
uses the range in the table, and $R$ is the remaining tail. For a
finite block $I$ put $v_I=2\sum_{i\in I}\gamma_i^{-2}$.

\begin{center}\small
\begin{tabular}{r r r r}
\hline
Indices & $v_I$ (rounded) & Integers sorted & Minimum gap${}^{*}$\\ \hline
1--500 & 0.04391240805 & \multicolumn{2}{l}{Best--Trudgian, Table 2}\\
501--6,000 & 0.00190153507 & 30,250,000 & 2,892\\
6,001--12,000 & 0.00016163380 & 17,997,000 & 236\\
12,001--18,000 & 0.00006286009 & 17,997,000 & 2,478\\ \hline
\end{tabular}

${}^{*}$In units of $10^{-14}$. The arithmetic uses unrounded enclosures.
\end{center}

For the first 500 ordinates, Best and Trudgian's Table 2 excludes every
nontrivial integer relation with coefficients of absolute value at most
$554,000$ (\href{https://arxiv.org/abs/1209.3843}{Math. Comp. 84 (2015), 2433--2446}).
In particular the 18th moment of $P$ equals the moment for independent
uniform phases. Indeed, a monomial in the expanded sine polynomial
has coefficient vector with $\ell^1$ norm at most 18. Only the constant
character integrates to one. For a single independent phase with
amplitude $2/\gamma$, the exponential moment series has coefficients
$1/(\gamma^{2k}(k!)^2)$ in degree $2k$ and zero in odd degree. It is
coefficientwise bounded by $\exp(z^2/\gamma^2)$ as a series in $z$.
Multiplication and degree extraction therefore give
\[
 \mathbb E[P^{2r}]\le (2r-1)!!\,v_P^r,\qquad r=9.
\]

For each middle block we prove that only pair cancellation contributes
to the fourth moment. We computed all 18,000 ordinates as certified balls
with FLINT, retaining the enclosing decimal cell of width $10^{-22}$.
Put $u_i=\lfloor10^{14}\gamma_i\rfloor$, computed from that cell.
The table records the minimum separation among positive differences
$u_j-u_i$ with $j>i$. In the first middle block we include in the same
sort all sums $u_j+u_i$ with $j\ge i$. Different positive differences
separated by more than three units exclude a nontrivial two-against-two
equality between ordinates, since such an equality would give two equal
positive differences after canceling common terms. A three-against-one
equality would give a sum equal to a positive difference. Their scaled
rounding errors are respectively in $[0,2]$ and $[-1,1]$, so the same
gap threshold suffices. In each of the last two blocks $3\min\gamma_i>
\max\gamma_i$ already excludes three-against-one. Four terms with a
common sign have positive absolute sum. Thus the fourth moment of
each $M_j$ equals its independent-phase value and is at most $3v_{M_j}^2$.
This certificate imposes no condition on relations between the blocks.

\paragraph{The part beyond the certificate.}
Brent, Platt and Trudgian give $\sum_{\gamma>0}m_\gamma/\gamma^2\le0.023105$
in equation (20) of \href{https://maths-people.anu.edu.au/brent/pd/rpb274-final.pdf}{J. Number Theory 238 (2022), 740--762}.
Platt and Trudgian verified simplicity to $H=3\cdot10^{12}$
(\href{https://doi.org/10.1112/blms.12460}{Bull. Lond. Math. Soc. 53 (2021)}).
Above $H$ a jump of the counting function is bounded by
$m_\gamma\le0.56\log\gamma$, by Corollary 1 of the 2022 paper.
Its Lemmas 5 and 6 bound the sums of $1/\gamma^2$ and
$\log(\gamma/2\pi)/\gamma^2$ beyond $H$, counting multiplicity.
With $D=\log H$ they give
\[
 \sum_{\gamma>H}\frac{m_\gamma(m_\gamma-1)}{\gamma^2}
 \le\frac{0.56}{2\pi H}\{D^2+(\log(2\pi)-1)D\}<2.524\cdot10^{-11}.
\]
The full variance is $2\sum m_\gamma^2/\gamma^2$ by orthogonality.
Subtracting the four finite block variances, with their certified
endpoints, yields $v_R=\|R\|_2^2<0.000171563042$.

\paragraph{A one-sided estimate without inter-block independence.}
Put $a_0=0.79$, $b_0=0.37$ and $M=\sum_j M_j$. Pointwise,
\[
 (F-a_0-b_0)_+\le(P-a_0)_+ +(M-b_0)_+ + R_+ .
\]
For $h>0$ and $r\ge2$, maximizing $(x-h)/x^r$ for $x\ge h$ gives
$(x-h)_+\le (r-1)^{r-1}(x_+)^r/(r^r h^{r-1})$.
With $r=9$ for $P$ and $r=2$ for $M$, symmetry and Minkowski give
\begin{align*}
 \|(P-a_0)_+\|_2 &\le
 \frac{(r-1)^{r-1}}{r^ra_0^{r-1}}\sqrt{(2r-1)!!\,v_P^r/2}<0.000923214,\\
 \|(M-b_0)_+\|_2 &\le
 \frac{\sqrt{3/2}}{4b_0}\left(\sum_{j=1}^3\sqrt{v_{M_j}}\right)^2<0.003415944,\\
 \|R_+\|_2&=\sqrt{v_R/2}<0.009261832.
\end{align*}
Chebyshev's inequality bounds the requested upper frequency by the
square of the sum of these norms divided by $(c-a_0-b_0)^2$.
To absorb the $o(1)$ in the explicit formula use threshold $c-\eta$,
discard a fixed initial interval at cost $O_\eta(1/L)$, then let
$\eta\downarrow0$. The unrounded computation gives an upper bound
below $0.00007942911$. Equality $E(x)=0$ occurs on a countable set,
since $E$ decreases strictly between consecutive prime powers.

\paragraph{An explicit lower bound for the rare sign.}
Let $G$ be the compact phase group above, with probability Haar measure,
and let $K_T$ be convolution along its one-parameter flow with the
nonnegative real-line Fej\'er kernel of integral one and Fourier multiplier
$(1-|\lambda|/T)_+$. Then $K_T$ preserves nonnegativity and constants,
is self-adjoint, and contracts the supremum norm. For $T=1420$ the first
$1{,}000$ ordinates lie below the cutoff, with the next above. Put
$d_i=2(1/\gamma_i-1/T)$ for $1\le i\le 1{,}000$. The smoothed sum
$K_TF$ is the finite polynomial $-\sum d_i\sin\theta_i$, where
$e^{i\theta_i}$ is coordinate $i$ on $G$.

Set $k=23$, and normalize $w(\theta)=(1-\sin\theta)^k$ to have
mean one on the circle. Direct integration gives
\[
  \max w=4^k/\binom{2k}{k}=:B,\qquad
  \int w(\theta)(-\sin\theta)\,\frac{d\theta}{2\pi}=k/(k+1).
\]
Indeed the circle mean before normalization is
$2^{-k}\binom{2k}{k}$, and the second identity follows by taking the
ratio of the successive means and subtracting one. Set
$W=\prod_{i=1}^{500}w(\theta_i)$. The finite Best--Trudgian certificate
applies coordinate by coordinate to every monomial of $W$ and of
$W\sin\theta_i$ for $i\le500$, since no exponent exceeds 24. Thus
$\int_G W=1$, while $0\le W\le B^{500}$. For the remaining 500 modes we
use only $|\sin\theta_i|\le1$. Therefore
\[
 \int_G W(K_TF-c)\ge
 \Delta :=\frac{k}{k+1}\sum_{i=1}^{500}d_i
     -\sum_{i=501}^{1000}d_i-c>0.00091075.
\]
Moving the convolution onto $W$ and using $0\le K_TW\le B^{500}$ gives
\[
 \Delta\le B^{500}\int_{F>c}(F-c)\le
 B^{500}\|F\|_2\,\mathbb P(F>c)^{1/2}.
\]
The arithmetic verifies $\Delta^2/(B^{1000}\|F\|_2^2)>10^{-937}$ using
the preceding upper bound for $\|F\|_2^2$.
To pass to uniform lower frequency, bounded Lipschitz test functions of
$F$ have Haar means uniformly on long intervals by the local mean-square
approximation above. Approximate the indicator of the open half-line
$(c,\infty)$ from below by such functions. The decaying explicit-formula
remainder also tends to zero in mean, uniformly on long intervals
starting at $a\ge1$. This proves the lower half of the theorem.

\clearpage
\section{Further research questions}
\paragraph{Nonabelian prime sums.}
For $L/\mathbb Q$ with group $S_n$ and $E=L^{\langle t\rangle}$, let
$r(p)$ be the number of fixed cosets of $\mathrm{Frob}_p$ on $S_n/\langle t\rangle$.
At unramified primes this equals $n!/2$ for the identity, $(n-2)!$ for
a transposition, and zero otherwise. Define
\[
 \theta_r(x)=\sum_{p\le x\ {\rm unramified}} r(p)\log p,\qquad
 R_r(x)=-\int_x^\infty (\theta_r(u)-u)\,\frac{du}{u^2}.
\]
The earlier reflection criterion would turn an upper bound
$R_r(x)\ll_{L,\eta}x^{-1/2+\eta}$ for every $\eta>0$ into RH for
$\zeta_L$ under Dirichlet GRH. Indeed the centered prime-power variant
of the transform has residue
$\operatorname{ord}_\rho\zeta_E/[\rho(1-\rho)]$, and dropping the higher
prime powers and ramified primes costs $O_L(x^{-1/2}\log x)$ after
centering. The source's geometric inequality has a rational congruence
condition and a quadratic splitting character as input.
Producing the corresponding geometric inequality with this nonabelian
Frobenius weight remains the substantive step. The sharp constants solve
the multiplicity problem for the groups treated here.

\paragraph{Finite groups.}
For $n\ge13$, which $S_n$ admit a positive monomial decomposition of the character
inequality in Section~\ref{sec:reflection} with coefficient $n-1$?
A family of identities uniform in $n$ would settle the earlier
optimization question in full. At a putative zero the standard character
already prevents improving the coefficient by character positivity alone.

\paragraph{Exceptional coefficient primes.}
Determine the gcd of \emph{all} positive integers satisfying the two Lean
conditions for general $N,e$. Theorem~\ref{th:gcd} evaluates it for a
specified certificate family, and the modulus-five theorem identifies an
instance where it equals the global gcd. An intrinsic answer could exploit
the Voronoi congruence without demanding that $2B_\chi$ divide $X$.
A separate question is whether the surviving primes actually support any
cocycles.

\paragraph{The boundary value in Suzuki's question.}
An atom-free limiting distribution at the particular threshold $c$
would give an ordinary logarithmic density. A finite list of ruled-out
resonances controls specified moments; it leaves possible relations with
arbitrarily high ordinates. A quantitative small-interval estimate at $c$
surviving the infinite-tail limit would address the missing point.
\clearpage
\appendix
\section{Coefficient collisions and higher-dimensional anticodes}
\label{sec:lattice}
\newcommand{\Z}{\mathbb Z}\newcommand{\R}{\mathbb R}

In the GRH project's construction of admissible data, the lemma
\texttt{exists\_three\_term\_free} chooses integer exponents $e$ modulo $o$.
For positive $a,b$ with $a+b\le L$, the congruence
$(a+b)e_2=ae_1+be_3$ forces $e_1=e_2=e_3$.
Its proof uses carry-free base expansions and strict convexity on a sphere.
The source identifies this step as a Behrend-type construction of known
provenance [1]. It suggests changing which coefficient relations must be
excluded, and asking how small an encoding group suffices.

Allow every equality between two sums of $h$ labels, with repetitions,
as a potential collision. A label set free of such collisions between distinct
multisets is a classical $B_h$ set. A coefficient difference
$z=(z_0,\ldots,z_n)$ has zero total and
$\tfrac12\sum_i|z_i|\le h$ after common labels cancel. Conversely, each
such difference comes from two $h$-term sums by padding both with the same
label. Thus the coefficient geometry is exactly that of the root lattice
\[
 A_n=\{z\in\Z^{n+1}:\textstyle\sum_i z_i=0\},\qquad
 d(z,w)=\tfrac12\textstyle\sum_i|z_i-w_i|.
\]
This classical correspondence [2] changes the GRH construction's modular
condition in a specific way. The consequences below are independent of the
analytic and automorphic theorems in that project.

\subsection{Exact higher-dimensional cases}
Write $a_n(D)$ for the largest cardinality of a subset of $A_n$ of diameter at
most $D$. Kova\v cevi\'c [3, \S7] conjectures that the balanced difference
of two discrete simplices attains $a_n(D)$ for every $n,D$. His theorem proves
this for $n=3$ and every $D$. The following are computer-assisted exact instances.
\begin{center}
\begin{tabular}{@{}r*{9}{r}@{}}\toprule
$D$ & 2 & 3 & 4 & 5 & 6 & 7 & 8 & 9 & 10\\\midrule
$a_4(D)$ & 21 & 55 & 131 & 255 & 471 & 780 & 1251 & 1875 & 2751\\
$a_5(D)$ & 31 & 96 & 271 & 606 & 1281 & & & &\\\bottomrule
\end{tabular}
\end{center}
For each displayed pair, every maximizing family is a translate of the balanced construction or its negative. The two $D=2$ entries also
follow from classical Singer difference sets for the prime powers $4$ and $5$,
through the code--anticode bound [2,3]. The cited theorem in rank three
does not cover either row for $D\ge3$. Priority for individual table
entries or for the small-group constants below has not been established.

For example, a family of $K$-multisets on five symbols in which every pair
shares at least $K-3$ elements, counting multiplicities, has size at most
$55$. Equality holds exactly when $K\ge6$.


\subsection{Proof of the table and the multiset interpretation}
For $x\in\Z^n$ put $p(x)=\sum_i\max(x_i,0)$ and $q(x)=p(-x)$.
The isometry $x\mapsto(x,-\sum x_i)$ identifies $d$ with the norm
$\max(p,q)$ of the difference. Put $P=\lceil D/2\rceil$,
$Q=\lfloor D/2\rfloor$. The candidate $S_n(P,Q)=\{x:p(x)\le P,\,q(x)\le Q\}$
has diameter at most $P+Q$, since $p(x-y)\le p(x)+q(y)$ and likewise
for $q$. Choosing the $j$ strictly positive coordinates and using
stars and bars on each sign gives its cardinality
\[
 M_n(D)=\sum_{j=0}^{\min(n,P)}
 {n\choose j}{P\choose j}{Q+n-j\choose n-j}.
\]

\textbf{Upper-bound algorithm.}
Let $F$ be an anticode in root coordinates, and put $m=n+1$.
Subtract each coordinate's minimum on $F$. All coordinates now lie
in $[0,D]$, and all sums equal some $k$. Reflection $x_i\mapsto D-x_i$
allows $0\le k\le\lfloor mD/2\rfloor$. Each possible ambient layer
$V_k$ is finite. For $x,y$ in it,
\[
 d(x,y)=\max_{J\subseteq\{1,\ldots,m\}}
 \sum_{i\in J}(x_i-y_i).
\]
Indeed, the sum of all positive coordinate differences is half their total
absolute value. Accordingly $F$ has diameter at most $D$ precisely when
each subset-sum has range at most $D$. Singletons already satisfy this,
and a cut and its complement have equal range.

For each remaining cut let $a,b$ be its minimum and maximum on the whole
layer. If $b-a>D$, try each integer interval $[s,s+D]$ for
$a\le s\le b-D$. Any feasible range is contained in one of these
intervals, by clipping its lower endpoint to $b-D$ when necessary.
If $b-a\le D$, the entire layer satisfies the cut. Intersect with a
chosen interval and recurse. Stop when the intersection has fewer than
$M_n(D)$ points. Every feasible family fits some branch at every depth.
Thus the absence of a larger intersection on all full branches proves the
upper bound. At equality the family must be the entire intersection.
For every leaf of size at least $M_n(D)$ the program checks its size
and compares the full set, after subtracting coordinate minima, with the
two normalized candidate shapes described below. This proves the equality
classification. The supplied \texttt{anticodes.py} performs these
operations with exact integer bit sets. Its outputs give the node counts.

\textbf{Multisets.} For two count vectors of sum $K$, intersection
cardinality is $K-d(x,y)$, so the upper bound holds on every layer.
For the lower bound take $x\in S_n(P,Q)$ and add final coordinate
$P-Q-\sum_i x_i$. All $m$ coordinates are at least $-Q$ and their sum is
$P-Q$. Each coordinate attains $-Q$, so the normalized shape is obtained by
adding $Q$ everywhere. It has coordinate sum $mQ+P-Q$. The negative
shape interchanges $P,Q$. Nonnegative translations give all larger sums.
The equality classification therefore gives the exact threshold
$K\ge mQ+P-Q$, since $P\ge Q$ and $m\ge2$. For $n=4,D=3$ this is $K\ge6$.

\textbf{Approximation in every fixed rank.}
Let $B=\{x\in\R^n:p(x),q(x)\le1\}$. Attaching unit cubes in the
$\Z^n$ coordinates to a family of size $a_n(D)$ gives a union $U$ of
that volume with $U-U\subseteq(D+n)B$. Brunn--Minkowski gives
$a_n(D)\le2^{-n}(D+n)^n\operatorname{vol}(B)$.
The formula for $M_n(D)$ gives leading coefficient
$2^{-n}{2n\choose n}/n!$. This equals $2^{-n}\operatorname{vol}(B)$.
Sum the orthant volumes $\binom nj/(j!(n-j)!)$ over $j$.
Consequently $a_n(D)=M_n(D)+O_n(D^{n-1})$. This elementary
approximation leaves the exact higher-dimensional conjecture open.

\subsection{The least order of a three-sum encoding with five labels}
Let $\phi(3,4)$ be the least order of a finite abelian group containing
a $B_3$ set of cardinality five, in the notation of [2].
The exact minimum is $64$. The minimum among cyclic groups is $65$.
Explicit witnesses in their respective groups are
\[
 \begin{array}{ll}
 (\Z/8\Z)^2 & (0,0),(0,1),(1,0),(4,6),(7,6),\\
 \Z/65\Z & 0,1,4,15,24.
 \end{array}
\]
Each witness gives $\binom{7}{3}=35$ distinct unordered three-label sums.

\textbf{Lower bound.} Translate a candidate so that $b_0=0$. The group
homomorphism $A_4\to G$ defined by $z\mapsto\sum_{i=0}^4z_i b_i$ has
no nonzero kernel element of distance at most $3$, by cancellation and
padding. It is injective on the $55$ points of $S_4(2,1)$ in root
coordinates. Thus $|G|\ge55$. For orders $55$ through $63$, the structure
theorem for finite abelian groups gives the cyclic group at every order
and just four additional groups with invariant factors
$(2,28),(2,2,14),(2,30),(3,21)$.

The check in \texttt{sidon.py} orders each group lexicographically and
enumerates all increasing four-label extensions of $\{0\}$. It retains
a partial set $T$ exactly when its unordered triple sums are distinct.
When adjoining a label $x$, the new sums are the elements of
$(x+2T)\cup(2x+T)\cup\{3x\}$, where $2T$ denotes the pairwise sumset.
There must be precisely $|2T|+|T|+1$ new values, all outside $3T$.
The cardinality test counts all new multisets correctly because the
partial $B_3$ property implies pairwise sums are distinct
by adding one fixed old label. Removing a label preserves the property,
so every valid set is reached. All thirteen groups below order $64$
in the specified range fail. The cyclic group of order $64$ also fails.
Both displayed witnesses pass an additional direct check over the $35$ multisets.

In particular, a full-rank sublattice of $A_4$ of minimum distance at
least four has index at least $64$. The five coordinate labels in its
quotient (translate a common coordinate vector to zero) have the $B_3$
property. A diameter-perfect lattice code at $D=3$ would have index
$a_4(3)=55$. This rules it out.

\textbf{Reproduction.}
From the directory \path{checks/lattice/} run
\begin{verbatim}
python anticodes.py 4 2 3 4 5 6 7 8 9 10
python anticodes.py 5 2 3 4 5 6
python sidon.py
\end{verbatim}

\paragraph{Sources.} [1] L. Sernau et al., \emph{GRH for Dirichlet
L-functions, in Lean 4}, readable edition,
\href{https://github.com/lsernau/grh-lean/blob/72fead102f663d5c7a407964a13dc0cf9d977207/grh/src/New/Reduction/Data/Existence.lean\#L57-L110}{\texttt{Existence.lean}, lines 57--110}.
[2] M. Kova\v cevi\'c and V.~Y.~F. Tan, \emph{Improved Bounds on
Sidon Sets via Lattice Packings of Simplices}, SIAM J. Discrete Math. 31
(2017), 2269--2278, \href{https://doi.org/10.1137/16M1099182}{DOI}.
[3] M. Kova\v cevi\'c, \emph{An Isodiametric Theorem and Lattice
Diameter-Perfect Codes in $A_3$}, 23 July 2026,
\href{https://arxiv.org/abs/2607.21037}{arXiv:2607.21037}, especially
\S\S2, 6--7.

\end{document}
